\documentclass[a4paper,14pt]{extarticle}
\usepackage{cmap}
\usepackage[T2A]{fontenc}
\usepackage[utf8]{inputenc}
\usepackage[russian]{babel}
\usepackage{geometry}
\geometry{left=30mm, right=15mm, top=20mm, bottom=20mm}
\usepackage{hyperref}
\usepackage{indentfirst}

\begin{document}

\begin{center}
    \textbf{\Large Оптимизации поиска базовых чанков в подходах SBC}
\end{center}

\vspace{1em}

\section*{Аннотация}

В данной учебной практике рассматривается актуальная задача оптимизации поиска базовых чанков в высоконагруженных системах дедупликации данных, работающих на основе оценки сходства (Similarity-Based Chunking, SBC). В таких системах для быстрого нахождения похожих блоков данных традиционно применяется инвертированный индекс. В нем хранятся уникальные признаки (features или sketches), извлеченные из каждого чанка. Основной проблемой является то, что при росте общего объема хранимых данных списки инвертированного индекса (posting lists) для часто встречающихся признаков становятся слишком длинными. В результате их полный обход требует чрезмерных вычислительных затрат, что делает процесс поиска узким местом всей системы, увеличивая задержки и потребление оперативной памяти.

Для преодоления этих ограничений в рамках работы исследуется и программно реализуется комплекс математически обоснованных поисковых эвристик для подсистемы \texttt{marline\_index}. Основной упор делается на интеллектуальную фильтрацию и отсечение заведомо нерелевантных кандидатов на самых ранних этапах поиска. В частности, внедряются такие оптимизационные механизмы, как отсев сверхчастотных признаков (\texttt{max\_df}), изменение порядка обхода индекса, начиная с самых редких признаков (\texttt{rare\_first}), и строгая фильтрация по порогу сходства Жаккара (\texttt{min\_jaccard}). Кроме того, разрабатывается система обратной связи, которая позволяет динамически штрафовать ложноположительные совпадения (\texttt{fp\_metric\_threshold}).

Особое внимание уделяется строгому математическому обоснованию каждого шага. Проводится тщательный статистический подбор пороговых значений и вычисляются безопасные границы для раннего отсечения. Это необходимо для того, чтобы многократное ускорение поиска (top-k) не приводило к случайному игнорированию действительно хороших базовых чанков и, как следствие, не снижало итоговый коэффициент дельта-компрессии.

Ожидаемым результатом работы является готовая, протестированная реализация набора эвристик на языке Rust, интегрированная в ядро \texttt{marline}. Дополнительно будет проведен сравнительный анализ производительности и бенчмаркинг на репрезентативных наборах данных, который наглядно продемонстрирует достигнутый баланс между скоростью поиска кандидатов и эффективностью итогового сжатия.

\end{document}
